\begin{abstract}
Retrieval-augmented generation systems increasingly flag contradictions among retrieved passages to
the generator. We ask what a flag must say, and how accurate its detector must be, to help against
misinformation. We vary detector and flag format on fixed passages from six benchmarks, with seven generators from three
families, testing hypotheses pre-registered (H1--H9 before any confirmatory data,
H10--H38 in nine later addenda). A neutral flag (``these passages conflict'') on gold pairs does not reliably
help. Its effect depends on answer policy and model family; pooled over \neuMlK\ settings, allowing
for shared questions, it is \neuMlMu\ points \neuMlCI\ (exploratory). Naming the false side helps, but on average no more than
an answer key naming the wrong answer (pooled difference \PoolTwentyThree\ points under our string metric; an outside-family judge favours the flag by \jdTwentyThreeJ). In a
randomised test, naming the false side raises accuracy by \hThirtyOneAMin\ to \hThirtyOneAMax\ points; naming the correct side lowers it by
\ahThirtyOneBMax\ to \ahThirtyOneBMin. Randomised flags on true pairs help only above a break-even direction
precision that rises with generator accuracy from \beMin\ to \beMax. Real detectors are selective, so this
threshold alone can misjudge them. Frozen before any flagged answer existed, the test's follow rates, with each cell's detections and no-flag accuracy,
predict flag effects in \frN\ new detector--generator cells on unseen questions (mean error \frMae\ points, slope \frSlope).
They overstate the harm on real health web pages, where
naming the false side still helps and a neutral flag does not. A detector's verdict, passed as an answer key built from benchmark labels, matches its flag's mean
accuracy under our metric. On genuine
conflicts, direction removes valid answers, and when conflicts are rare a real detector's false flags on clean
questions can outweigh its gains. We release the harness, pre-registrations and every run.
\end{abstract}
